\section{Discussion}\label{sec:discussion}
\subsection{Answers to the research questions}
The first objective was to assess whether ML improves on seasonal rules under chronological evaluation. Several default models improve on Naive-24h in November, yet their ordering changes across temporal folds. LR is the strongest default by mean CV RMSE, while tuning selects absolute SVR. The selected h24 SVR reduces December RMSE by \num{50.71}\% against the daily naive reference and also improves on the weekly reference. Thus improvement is supported within the observed evaluation setting, rather than as an unconditional ranking of algorithms.

The feature experiments show that daily history is the most effective sequential addition tested for RF/XGB. Calendar variables help represent recurring demand differences, but weather and calendar inputs without load history remain poor substitutes for a recent measured level. Weekly and rolling summaries worsen results under D-selected parameters. The conclusion is to match feature construction to validation evidence; the study does not establish that weekly information is intrinsically useless or that the same subsets would win after independent retuning.

The residual formulation improves tuned development mean RMSE for four of five families, helping trees predict corrections around an available recent load. Absolute SVR remains slightly better in development for both settings. The different h24 test ordering shows why the choice must be retained separately from the observed test winner. The h1 configuration improves mean CV RMSE for each family representative, but its feature set, residual base and tuning differ from h24. The study supports this configuration-level finding without claiming a controlled lead-time response curve.

\subsection{Accuracy, complexity and practical cost}
\Cref{tab:runtime} gives the recorded h24 refit and prediction durations. The selected SVR takes approximately \num{3.686} s to fit and \num{0.137} s to predict 720 rows. Residual RF and residual XGB have shorter recorded fits, approximately \num{0.107} and \num{0.146} s, with test RMSEs \num{1761.9} and \num{1719.2} kW. Their development RMSEs are also close to SVR's. These alternatives may be relevant if latency or tree-based explanations matter, but the experiment does not apply a formal multi-objective selection rule.
\begin{table}[tbp]\centering\small
\caption{Recorded h24 refit and whole-test prediction durations for selected representative models (T08). Hardware and thread configuration are not recorded in this table; timings are descriptive.}\label{tab:runtime}
\setlength{\tabcolsep}{4pt}
\begin{tabular}{lrrr}
\toprule
Model & Fit (s) & Predict 720 rows (s) & Test RMSE (kW) \\
\midrule
LR $\Delta$-tuned & \num{0.003} & \num{0.001} & \num{2217.3} \\
DT $\Delta$-tuned & \num{0.015} & \num{0.001} & \num{2005.2} \\
RF $\Delta$-tuned & \num{0.107} & \num{0.026} & \num{1761.9} \\
SVR-tuned & \num{3.686} & \num{0.137} & \num{1624.6} \\
XGB $\Delta$-tuned & \num{0.146} & \num{0.003} & \num{1719.2} \\
\bottomrule\end{tabular}
\end{table}


The timings are single recorded runs without hardware provenance, repeated benchmarking or a consistent deployment service. They do not establish production throughput or scaling to larger datasets. Likewise, the reported number of search fits describes computational work, not an energy or monetary cost measurement. The uncomplicated seasonal and linear references remain useful as cheap diagnostics and potential fallback models, although a fallback policy was not implemented.

\subsection{Validation reliability and generalization}
Expanding folds expose summer errors hidden by a November-only evaluation. This is a concrete benefit of temporal validation, not a guarantee of future generalization. Fold scores combine seasonal difficulty with increasing training sizes, and reuse for feature choices and tuning creates selection bias. The narrow spread among leading tuned means is much smaller than their temporal dispersion; neither a significant winner nor statistical equivalence has been demonstrated.

The final evaluation covers a single 30-day period in one city/year. Its lower error than average development CV cannot represent all seasons. The early-December holiday failure also indicates that broad special-day categories may pool events with different demand responses. Sparse holiday occurrences prevent a reliable event-specific performance estimate. No zone-level forecasting results, multi-year transfer study, forecast intervals or external-city validation are available.

The measurement context further limits deployment claims. Observed target-hour weather overstates available information for advance forecasting unless replaced by an issuance-time forecast. Historical-load delivery latency and clock conventions are unspecified. Test reruns and descriptive exposure to December prevent certifying a fully hidden one-shot evaluation, even though the selection expression uses development scores. These are material boundaries on what the current evidence demonstrates.

\subsection{Relationship to verified prior work}
The literature in \cref{tab:literature} provides methodological context. Salam and El Hibaoui motivate conventional regression comparisons\cite{salam2018}; da Silva and Meneses provide a chronological deep-learning comparison as described in the local review\cite{silva2023}; shift-based and lag-enriched approaches emphasize historical demand\cite{islam2025,alsalem2025}; and the Darts study illustrates comparison across model families and resolutions\cite{biswas2026}. The present findings are consistent with the usefulness of history and careful validation, but those studies do not share a verified common target, lead, covariate-availability protocol and test period with this implementation.

For example, a zone-level ten-minute MAE cannot be ranked directly against an hourly three-zone total MAE. Taking the mean of three zones rather than their sum changes numerical scale, but correcting that scale alone does not align horizons or evaluation distributions. The report therefore retains traceable study descriptions and unavailable entries instead of constructing a cross-study leaderboard. F20 is not reproduced because its heterogeneous bar comparisons could encourage precisely that interpretation. No state-of-the-art conclusion follows.

\subsection{Implications and remaining evidence gaps}
For a practical forecasting workflow, the strongest supported lessons are to retain seasonal baselines, define the issuance schedule, audit every predictor's availability, and keep model-selection scores separate from final test observations. Interpretability should identify which estimator and output it explains: the saved Tree SHAP plots describe residual RF, not the selected SVR. Permutation experiments on correlated calendar/history columns should be treated as model diagnostics rather than causal importance.

Further claims require evidence absent from this snapshot: weather forecasts as issued, measured data latency, independently reserved seasonal/yearly tests, per-zone experiments, repeated computational benchmarks, uncertainty estimates and archived h1 forecast series. Some historical side experiments are described in documentation without their full fold/search artifacts. They can explain development rationale but cannot substitute for reproducible numerical evidence. All these gaps are recorded in the research audit rather than filled by new experiments for this report.
\evidence{T04--T10; T12/T12b; related-work evidence in \code{research_audit.md}; \code{docs/result.md}, Section 11; recorded runtimes in T08.}
